% ==========================================
% 1. TERMINI DI GLOSSARIO (Le definizioni estese)
% ==========================================

\newglossaryentry{rpg_ext}{
    name={Role-Playing Game (RPG)},
    description={Un gioco dove i giocatori assumono il ruolo di uno o più personaggi e tramite la conversazione e lo scambio dialettico creano uno spazio immaginato, dove avvengono fatti ed eventi fittizi, in un'ambientazione narrativa che può ispirarsi a un romanzo, a un film o a un'altra fonte creativa, storica, realistica o di pura invenzione. Esiste anche il termine in italiano GDR (Gioco di Ruolo), ma è poco usato}
}

\newglossaryentry{mmorpg_ext}{
    name={Massive Multiplayer Online Role-Playing Game (MMORPG)},
    description={Un videogioco di ruolo (per \textit{computer} o \textit{console}) che viene svolto giocato da più persone contemporaneamente attraverso Internet. Negli MMORPG, migliaiadi giocatori possono giocare insieme ed interagire tra di loro, interpretando personaggi personalizzabili che si evolvono (acquisendo abilità, guadagnando valuta \textit{in-game}, ottenendo oggetti ed equipaggiamenti, ecc.) insieme al mondo persistente che li circonda e in cui si trovano}
}

\newglossaryentry{api_ext}{
    name={Application Programming Interface (API)},
    description={Un insieme di procedure (in genere raggruppate per strumenti specifici) atte a consentire la comunicazione tra diversi \textit{computer} o tra diversi \textit{software}; spesso tale termine designa le librerie \textit{software} di un linguaggio di programmazione, sebbene più propriamente le API siano il metodo con cui le librerie vengono usate per sopperire ad uno specifico problema di scambio di informazioni}
}

\newglossaryentry{sdk_ext}{
    name={Software Development Kit (SDK)},
    description={Un insieme di strumenti per lo sviluppo e la documentazione di \textit{software}, che possono variare anche in base alla specializazzione nel caso di domini complessi. Alcuni strumenti fondamentali presenti sono: un compilatore, un }
}

\newglossaryentry{obb_ext}{
    name={Opaque Binary Blob (OBB)},
    description={\textit{File} di espansione binario \textit{opaque} impiegato principalmente nell'ecosistema Android per distribuire risorse e contenuti multimediali di grandi dimensioni (come \textit{asset} grafici, modelli 3D o audio) separatamente dal pacchetto principale dell'applicazione (\textit{file APK}).}
}

% ==========================================
% 2. ACRONIMI COLLEGATI (Popolano entrambe le liste)
% ==========================================
% Nota il \glsadd nascosto alla fine del nome per esteso

\newacronym{rpg}{RPG}{Role-Playing Game\glsadd{rpg_ext}}
\newacronym{mmorpg}{MMORPG}{Massive Multiplayer Online Role-Playing Game\glsadd{mmorpg_ext}}
\newacronym{api}{API}{Trading Card Game\glsadd{api_ext}}
\newacronym{sdk}{SDK}{Software Development Kit\glsadd{sdk_ext}}
\newacronym{obb}{OBB}{Opaque Binary Blob\glsadd{obb_ext}}

% ==========================================
% 3. ACRONIMI SEMPLICI (Vanno solo nella lista acronimi)
% ==========================================

\newacronym{tcg}{TCG}{Traiding Card Game}

% ==========================================
% 4. TERMINI DI GLOSSARIO STANDARD
% ==========================================

\newglossaryentry{dungeon crawler}{
    name={Dungeon Crawler},
    description={Un gioco derivante direttamente dagli \textit{RPG}, con la presenza di labirinti noti come dungeon e sono caratterizzati da un'ambientazione fantastica con la presenza di creature mitologiche e incantesimi magici}
}

\newglossaryentry{game engine}{
    name={Game Engine},
    description={Un \textit{software} usate prevalentemente per lo sviluppo di videogiochi e la loro distribuzione su diverse piattaforme. Comunemente offrono funzionalità come il rendering 2D e 3D, un motore di simulaizione fisico, animazioni e altro ancora}
}

\newglossaryentry{scriptable object}{
    name={Scriptable Object},
    description={Gli Scriptable Objects sono contenitori di dati in Unity che ti permettono di memorizzare e gestire i dati. Possono anche contenere logica e sono spesso utilizzati per centralizzare i dati da utilizzare su più oggetti in un progetto}
}

\newglossaryentry{publish-subscribe}{
    name={Publish-Subscribe},
    description={Schema di messaggistica in cui i mittenti dei messaggi, chiamati \textit{publisher}, categorizzano i messaggi in classi e li inviano senza dover sapere quali componenti li riceveranno}
}

\newglossaryentry{Issue Tracking System}{
    name={Issue Tracking System},
    description={Un \textit{software} utilizzato per organizzare e gestire "problematiche" di ampio genere, dalla creazione alla risoluzione, solitamente con proprietari, priorità, stati, etichette e flussi di lavoro}
}

\newglossaryentry{build}{
    name={Build},
    description={Termine utilizzato, in ambito informatico, per indicare il processo di trasformazione del codice sorgente in un artefatto eseguibile}
}

\newglossaryentry{run}{
    name={Run},
    description={Nel linguaggio videoludico viene usata in modo aleatori per indicare una partita oppure un tentativo specifico nei giochi in cui tutto il contenuto si svolge in un'unica sessione}
}

\newglossaryentry{tutorial}{
    name={Tutorial},
    description={Fase introduttiva del gioco, in cui vengono insegnate le meccaniche base al giocatore}
}

\newglossaryentry{mana}{
    name={Mana},
    description={Termine di origine melanesiana che indica generalmente una forza sovrannaturale. In contesti di storie fantasy e nei vidoegiochi rappresenta l'energia magica con cui si possono eseguire gli incatesimi}
}